\notechapter{N-20220913}{Divisibility, Congruences, and the Habit of Proving Necessity}

\section{From guessing to proving}

The note opens with a useful warning: in number theory, guessing examples is not enough.  One must separate necessity and sufficiency.  If a condition is necessary, it must follow from the assumption.  If it is sufficient, it must imply the desired conclusion.  Many false solutions prove only one direction.

The first example asks for a prime \(p\) of the form
\[
  p=x(x+1)-4,\qquad x\in\mathbb N_+.
\]
Since \(x(x+1)\) is always even,
\[
  p=x(x+1)-4
\]
is even.  If \(p\) is prime, then the only even prime is \(2\).  Hence
\[
  p=2,\qquad x(x+1)=6,
\]
so
\[
  x=2.
\]
This is a clean example of using parity as a necessary condition before solving.

\section{Divisibility}

For integers \(a,b\), with \(a\neq0\), we write
\[
  a\mid b
\]
if there exists \(q\in\mathbb Z\) such that
\[
  b=aq.
\]
The standard properties are:
\begin{enumerate}[(i)]
  \item \(a\mid b\Longleftrightarrow -a\mid b\Longleftrightarrow a\mid -b\Longleftrightarrow |a|\mid |b|\);
  \item if \(a\mid b\) and \(b\mid c\), then \(a\mid c\);
  \item if \(a\mid b\) and \(a\mid c\), then \(a\mid bx+cy\) for all integers \(x,y\);
  \item if \(m\neq0\), then \(a\mid b\Longleftrightarrow ma\mid mb\);
  \item if \(a\mid b\) and \(b\mid a\), then \(b=\pm a\);
  \item if \(b\neq0\) and \(a\mid b\), then \(|a|\le |b|\).

\end{enumerate}

Property (iii) is often the most useful because it combines divisibility relations linearly.

\section{Example: reducing the dividend}

Suppose
\[
  x^2+1\mid x^3-x,\qquad x\in\mathbb Z_+.
\]
Modulo \(x^2+1\), we have
\[
  x^2\equiv -1,
\]
so
\[
  x^3-x=x(x^2)-x\equiv -x-x=-2x.
\]
Thus
\[
  x^2+1\mid 2x.
\]
For \(x>1\),
\[
  x^2+1>2x,
\]
so divisibility is impossible unless \(2x=0\), which cannot happen for positive \(x\).  Checking \(x=1\), we get
\[
  1^3-1=0,
\]
so \(x=1\) works.

The strategy is to reduce the large expression using the divisor itself and then use size estimates to force a small list of cases.

\section{Example: polynomial divisibility in one variable}

For divisibility such as
\[
  x^2+x+1\mid x^4-2,
\]
one uses
\[
  x^2+x+1=0\quad\Longrightarrow\quad x^3-1=(x-1)(x^2+x+1)\equiv0.
\]
Thus \(x^3\equiv1\), and
\[
  x^4-2\equiv x-2.
\]
So the divisor must divide \(x-2\).  This leaves only very small candidates, because for large \(|x|\), the quadratic divisor is too large in absolute value to divide the linear remainder.

The general method is Euclidean division in disguise: replace the dividend by its remainder modulo the divisor.

\section{Congruences}

Congruence notation compresses divisibility:
\[
  a\equiv b\pmod m
\]
means
\[
  m\mid a-b.
\]
This turns many divisibility problems into arithmetic with remainders.  Congruences can be added, subtracted, and multiplied:
\[
  a\equiv b\pmod m,\quad c\equiv d\pmod m
  \Longrightarrow
  a+c\equiv b+d,\quad ac\equiv bd\pmod m.
\]

\section{A large power sum}

The last problem asks to prove a divisibility statement involving
\[
  1^{1983}+2^{1983}+\cdots+1983^{1983}.
\]
The intended thought is pairing and congruence.  When the exponent is odd,
\[
  (m-k)^{1983}\equiv -k^{1983}\pmod m.
\]
So terms can be paired to cancel modulo a suitable modulus.  To prove divisibility by
\[
  1+2+\cdots+1983=\frac{1983\cdot1984}{2},
\]
one decomposes the modulus into factors and proves divisibility factor by factor, using pairing and residue classes.

For large power sums, first look for the symmetry \(k\leftrightarrow m-k\), then reduce the problem to congruences.

\section{A wider interpretation}

Divisibility is the first place where algebraic thinking becomes economical.  Instead of expanding huge expressions, we work modulo the divisor and keep only remainders.  This is exactly the same principle behind quotient rings such as
\[
  \mathbb Z/m\mathbb Z
\]
and polynomial quotients
\[
  F[x]/(f).
\]
These small examples are elementary instances of modular algebra.

\section{Divisibility as ideal containment}

In \(\mathbb Z\), the statement \(a\mid b\) is equivalent to
\[
  b\in(a),
\]
where \((a)=a\mathbb Z\) is the ideal generated by \(a\).  Divisibility reverses inclusion of principal ideals:
\[
  a\mid b\quad\Longleftrightarrow\quad (b)\subseteq(a).
\]
This explains why linear combinations are central.  If \(a\mid b\) and \(a\mid c\), then every integer linear combination \(bx+cy\) lies in \((a)\).

\section{GCD and LCM as ideal operations}

For integers \(a,b\), Bezout's theorem says
\[
  (a)+(b)=(\gcd(a,b)).
\]
The ideal generated by all linear combinations of \(a\) and \(b\) is generated by their greatest common divisor.  Similarly,
\[
  (a)\cap(b)=(\operatorname{lcm}(a,b)).
\]
Thus gcd and lcm are not just computational devices; they are sum and intersection of ideals in \(\mathbb Z\).

\section{Congruences as quotient rings}

The congruence
\[
  a\equiv b\pmod n
\]
means that \(a\) and \(b\) define the same element of the quotient ring
\[
  \mathbb Z/n\mathbb Z.
\]
The projection
\[
  \pi:\mathbb Z\to\mathbb Z/n\mathbb Z
\]
forgets the exact integer and remembers only its residue class.  Its kernel is \(n\mathbb Z\), so this is the first isomorphism theorem in its simplest form:
\[
  \mathbb Z/n\mathbb Z\cong \mathbb Z/\ker\pi.
\]
Modular arithmetic is therefore quotient algebra, not merely a remainder shortcut.

\section{Units and Bezout inverses}

An integer \(a\) has a multiplicative inverse modulo \(n\) exactly when
\[
  \gcd(a,n)=1.
\]
By Bezout, there exist integers \(x,y\) such that
\[
  ax+ny=1.
\]
Reducing modulo \(n\) gives
\[
  ax\equiv1\pmod n,
\]
so \(x\) is the inverse of \(a\) modulo \(n\).  This is the concrete connection between gcd computations and units in \(\mathbb Z/n\mathbb Z\).

\section{Valuations as prime coordinates}

For nonzero integers, \(v_p(n)\) denotes the exponent of the prime \(p\) in \(n\).  Divisibility becomes coordinatewise comparison:
\[
  a\mid b\quad\Longleftrightarrow\quad
  v_p(a)\le v_p(b)\text{ for all primes }p.
\]
Gcd and lcm become coordinatewise minimum and maximum:
\[
  v_p(\gcd(a,b))=\min(v_p(a),v_p(b)),\qquad
  v_p(\operatorname{lcm}(a,b))=\max(v_p(a),v_p(b)).
\]
This turns divisibility into geometry on the exponent lattice of prime factorizations.

\section{Chinese remainder theorem as local-to-global}

\begin{theorem}[Chinese remainder theorem]
If \(m_1,\ldots,m_r\) are pairwise coprime positive integers, then
\[
  \mathbb Z/(m_1\cdots m_r)\mathbb Z\cong
  \prod_{i=1}^r \mathbb Z/m_i\mathbb Z.
\]
\end{theorem}


If
\[
  n=\prod_i p_i^{e_i},
\]
then the Chinese remainder theorem gives
\[
  \mathbb Z/n\mathbb Z\cong
  \prod_i \mathbb Z/p_i^{e_i}\mathbb Z.
\]
So a congruence problem modulo \(n\) can often be solved prime-power by prime-power and then recombined.

This is a first local-to-global principle: understand arithmetic at each prime separately, then assemble the global answer.

\section{Polynomial divisibility as the same story}

Over a field \(k\), the ring \(k[x]\) has Euclidean division:
\[
  f=qg+r,\qquad \deg r<\deg g.
\]
Therefore the methods of integer divisibility have polynomial analogues.  In particular,
\[
  x-a\mid f(x)\quad\Longleftrightarrow\quad f(a)=0.
\]
Working modulo a polynomial \(g(x)\) means working in the quotient ring \(k[x]/(g)\).  The earlier examples that reduce \(x^4-2\) modulo \(x^2+x+1\) are exactly quotient-ring computations in elementary disguise.

\section{Divisibility as quotient algebra}

Divisibility problems become clearer when the divisor is treated as a structure rather than an obstacle.  Ideals organize divisibility, quotient rings organize congruence, valuations localize divisibility at primes, and Euclidean division transfers the same ideas to polynomials.  The elementary habit "reduce by the divisor" is the first step toward quotient algebra.
